\section{Measure work after the arithmetic agrees}
Count simulated instructions over a stated function boundary, including its normal
saves, nested calls and return. Record the base, exponent, modulus, scan direction and
whether the input reduction took its direct-copy path. These choices affect work.
A sparse exponent such as 17 has only two set bits; an all-ones 128-bit exponent
selects a multiply at every bit. A fixed 128-bit scan also visits leading zero bits
that a significant-bit scan can skip.

Predict the number of selected multiplies and squarings before measuring. For a
nonzero exponent with bit length \(L\), the right-to-left loop shown earlier performs
\(L-1\) squarings and one result multiply for every set exponent bit. Its zero-exponent
case performs neither. The nested cost still depends on residue values and multiplier
bits in this teaching implementation. Instruction counts describe simulator work;
they do not measure wall-clock latency, hardware cycles or constant-time behavior.

The current complete cumulative implementations give the following counts for base 65 and modulus
\(2^{128}-159\). The first scans only the significant remaining bits from right to
left. The second scans all 128 bits from left to right, including leading zeros.
Both include the nested status guards used by the construction.
The measured boundary starts at the function entry and ends at its caller's return
label; caller setup, assembly and session construction are excluded.
\tagpdfsetup{table/header-rows={1}}
\begin{center}
\begin{tabular}{lrr}
Exponent & Right-to-left & Fixed left-to-right \\
\hline
0 & 130 & 26,624 \\
1 & 1,053 & 27,533 \\
17 & 27,606 & 40,076 \\
2753 & 253,857 & 198,056 \\
\(2^{127}\) & 2,790,405 & 2,771,162 \\
\(2^{128}-1\) & 5,573,877 & 2,948,849 \\
\end{tabular}
\end{center}
The fixed scan spends work even on exponent zero. It nevertheless uses fewer simulated
instructions for some larger exponents in this comparison. Its intermediate residues
and the operands supplied to nested multiplication differ. These totals combine those
value-dependent costs with scan control and storage; they do not isolate one instruction
or establish that either scan direction is always faster. Both implementations produce
the same checked output. Exact totals are observations of these prototypes, not grading
requirements or security properties.

First investigate repeated argument setup, copies and reductions of already bounded
residues. A private helper whose preconditions explicitly require residues can avoid
some public-entry work, provided callers establish those bounds. Keep the public
interfaces unchanged and run the same value, aliasing, preservation and address checks
after a change. Do not optimize by dropping a high limb or assuming that every modulus
is below \(2^{127}\).

As further exercises, derive a 256-bit bitwise reducer for the full product, or extend
the limb count beyond four. State the larger buffer sizes and loop bounds before
changing assembly. Compare its work with add-and-double modular multiplication. These
extensions require new checks; the four-limb cases cannot establish an arbitrary-length
integer library. More advanced reductions also need their own modulus preconditions
and derivations.
